\documentclass[11pt]{article}

\usepackage[T1]{fontenc}
\usepackage[margin=0.68in]{geometry}

\pagestyle{empty}
\setlength{\parindent}{0pt}
\setlength{\parskip}{0.32em}

\begin{document}

\begin{center}
    {\Large \textbf{Earnings Surprise, Firm Valuation, and Market Efficiency:}\\[-0.15em]
    \textbf{Evidence from China}}\\[0.2em]
    {\small Team Project Proposal for Corporate Finance (FIN803)}\\[0.1em]
    {\footnotesize Team Members: Yingqi Tan, Zhesheng Xu, Hao Yang}
\end{center}

\small

\noindent\textbf{Project Title.} Earnings Surprise, Firm Valuation, and Market Efficiency: Evidence from China.

\noindent\textbf{Background and Motivation.} Earnings-related disclosures are important information events because they update investors' expectations about future profitability and firm value. In this project, the main focus is not whether raw earnings news always produces long-run drift, but whether a cleaner and more defensible surprise measure predicts short-run abnormal return after public guidance announcements in the China A-share market. This question is especially relevant in China, where investor composition, information processing frictions, and institutional features may produce limited short-run delays in price adjustment even if strong long-run PEAD is not robustly supported.

\noindent\textbf{Research Question.} The central question is whether cleaner standardized guidance surprise in China predicts post-announcement abnormal return over a short horizon. More broadly, we ask whether earnings-related public information is incorporated into prices immediately, or whether there is limited short-run valuation adjustment after guidance news.

\noindent\textbf{Relevance to FIN803 Corporate Finance.} This topic is closely aligned with the FIN803 syllabus. First, it relates to \textbf{stock valuation} because guidance updates expectations about future cash flows and therefore affects intrinsic value. Second, it connects to \textbf{capital markets and the pricing of risk} by separating normal expected return from event-related abnormal return. Third, it uses ideas from \textbf{CAPM and cost of capital} to benchmark returns and evaluate whether post-announcement return patterns reflect risk or delayed information incorporation. The final interpretation is intentionally conservative: the project studies valuation adjustment around public information, not an aggressive inefficiency or trading-strategy claim.

\noindent\textbf{Methodology.} We study Chinese listed firms and earnings-guidance announcements. The preferred signal is a standardized guidance surprise measure, denoted \(ES_{std}\), designed to reduce noise relative to the earlier proxy-based surprise definition. The main dependent variable is short-window abnormal return over five trading days after the announcement, \(CAR5\). Abnormal return is defined as stock return minus market return. The preferred regression relates \(CAR5\) to \(ES_{std}\), while controlling for firm size and beta, and including industry fixed effects, year-quarter fixed effects, and firm-clustered standard errors. This specification is more defensible than a long-window \(CAR60\) design because it is less contaminated by unrelated subsequent news and better isolates the market's short-run reaction to guidance information.

\noindent\textbf{Main Finding and Contribution.} The updated evidence suggests that cleaner standardized guidance surprise predicts a \textit{modest but statistically significant} short-run abnormal return in China A-shares. The contribution of the project is therefore not to claim strong long-run PEAD, but to show that better-measured guidance surprise is associated with limited short-run valuation adjustment after public disclosure. From a Corporate Finance perspective, this means that earnings-related information does matter for firm value, but the magnitude of delayed adjustment is economically limited and should be framed cautiously. Overall, the project provides a more academically defensible interpretation of how accounting-related information enters prices in an important emerging market.

\end{document}
